\section{FlexOS}

\subsection{Nombre del proyecto o sistema operativo}

El sistema operativo examinado es "Redox OS". De acuerdo con \citep{redox_docs}, su denominación deriva de la palabra “Redox”, que es la abreviatura de “\textit{reduction–oxidation}”, un término de química que representa el intercambio equilibrado de electrones; esto refleja la meta del proyecto: alcanzar un sistema equilibrado en seguridad, rendimiento y simplicidad, a través de un diseño contemporáneo y fiable, redactado totalmente en Rust.

\subsection{Enlace al repositorio y/o documentación oficial}
\begin{itemize}
    \item \textbf{Enlace al repositorio oficial:} https://github.com/redox-os
    \item \textbf{Enlace al libro oficial:} https://doc.redox-os.org/book/
\end{itemize}

\subsection{Objetivo del proyecto}

\subsection{Lenguaje(s) de implementación}


\subsection{Arquitectura del sistema }



\subsection{Componentes implementados }



\subsubsection*{Gestión de memoria}


\subsubsection*{Carga dinámica de módulos (módulos / células)}


\subsubsection*{Scheduling / multitarea}



\subsubsection*{Subsistemas de E/S y drivers}



\subsubsection*{Recuperación ante fallos / actualización en caliente}


\subsubsection*{Comunicación entre módulos}


\subsection{Herramientas utilizadas (compiladores, emuladores, etc.)}



\subsection{Nivel de complejidad y accesibilidad para estudiantes}